\section{Gossip protocols}
Gossip protocols (also known as epidemic protocols) are a model to spread a piece of information among a large number of processes with a dynamic connection topology. They exhibit a very robust and scalable features even in the presence of a high rate of link failures. They have been studied theoretically and their analysis is built on mathematical foundations. The survey \cite{EUG04} provides an introduction to the field and delineates some key issues while designing a gossip protocols.

\paragraph{Typical application} The use of epidmics algorithms has been explored in applications such as database replication, failure detection, data aggregation, resource discovering and monitoring. We will make use of epidemic protocols in order to dinamically build a network topology which reflects the semantic clustering of information.

\paragraph{Model}The principle underlying this information dissemination techniques mimics the spread of epidemics. Another close analogy is with the spread of a rumour among humans via gossiping.

A process that wishes to disseminate a new piece of information to the population, does not need to know all the partecipants, which membership could have high dynamics, but only a small subset of other peer processes.

Every process buffers every message (information unit) it receives up to a certain buffer capacity $b$. It forwards that message a limited number $t$ of times to a randomly chosen subset of processes of limited size $f$ (i.e. the fanout of the dissemination). The reliability of information delivery will depend both on these values as well as on the size $n$ of the whole population.

\section{Locality of interest principle} Locality of interest principle, as defined by \cite{SMZ03} posit that if a peer has a particular piece of content that one is interested in, it is very likely that it will have other items that one is interested in as well. They demonstrate this assumption by examining traces of content distribution applications. such p2p file sharing systems. It has been observed that p2p file sharing traces exhibit locality of interests \cite{SMZ03}.

\section{Topology}
Characteristics of the communication topology play a key role in determining the performance of many functions including search, routing, monitoring, control and information dissemination.

Network topology can also be used to define solutions for problems such as clustering nodes based on similarity or pairing search requests with the location of the information.

Topology has to support the application that is being implemented. For example in the case of routing or search. the routing tables have to define a low diameter topology that also includes some information hints on how to get closer to the destination



\subsection{Topology construction problem}

\paragraph{Problem statement} We consider a set of processes distributed over a network. Each process is reachable by using a network address and a port, that is necessary and sufficient for sending it a message. Each node knows about $c$ other nodes of the networks through a partial view which is a set of node descriptors. The local knowledge of each node about the network, i.e. their partial views define the links of the overlay network topology.

In addition each node has a profile as part of its descriptor.

The input of the problem is a set of $N$ nodes, the view size $c$ and a ranking fuction $R(P_i, P_j)$ defined over the node profiles which defines (at least) a partial ordering of the nodes $y_1, \ldots, y_m$.

\begin{equation}
 d: P \times P \mapsto \mathbb{R}
\end{equation}

The goal is to build the views of the nodes such that for each node $y_i$ its view $\mbox{view}_i$ contains exactly the first $c$ elements in the order as defined by the scoring fuction $R$ out of the whole network.

\paragraph{Ranking functions} The topology of the network is then defined by the choice of the ranking function which ranks node according to the defined order. Ir can be based on an approximate property of components including geographical location, semantic proximity, available bandwidth or simply on an abstract, pre-defined topology over an ID-space like a ring or a torus.

One way of generating ranking functions is through the distance function $d(P_i,P_j)$ where the profiles $P_i, P_j$ of the nodes are elements in a metric space. In such case the ranking function can simply order the given set according to the distance from elements.

Because we are interested in clustering node which are close in their semantics we should chose a proper ranking function defined over the node's profile which model this need.

A metric space is such that for any element in the set it is tr.
\begin{equation}
 d(P_x,P_z) \leq d(P_x,P_y) + d(P_y,P_z)
\end{equation}

It is desirable to chose a ranking function which holds the triangle inequality because it eases the use of epidemic algorithms to manage topologies.

\paragraph{Reactive vs. proactive approaches} There are two strategies on exploiting network topologies for improve the content searching.

Reactive approach make usage of the underlying search mechanism and infer semantic relations based on the query placed and the corresponding replies received. The network topology is changed each time a user pose a query to the system \cite{VKM+04, SMZ03}. They generally rely on heuristics to decide which peers that served a node recently are likely to be useful again for future queries. One of the hidden assumptions those solutions make is that the network is static both in terms of node joining and leaving and changes in nodes' profiles. \cite{VKM+04} examines different policies in order to modify the network topology to improve search performance.

Proactive approach builds a semantic overlay topology by making use of a peer profile and a ranking function which define an order over the peers. The topology construction is done proactively in a competely implicit way, that is decoupled to the execution of queries. The peer profile could capture the user's search trends through analysis of query posed, the results, his file cache.

The reactive approach minimize the resource usage, but leads to low hit ratio. Proactive approach his more intrusive but it can lead in high quality semantic overlay topologies, based on the choice of a proper ranking function and a peer profile.

\paragraph{Concernings on managing topologies by using an epidemic approach} There are two sides to the construction of a topology over a set of nodes distributed in a network.

Exploiting the triangle inequality property of the ranking function $R(P_i,P_j)$ should quickly lead to high quality local views. For instance if $Q$ is in the local view of $P$ and $R$ is in the local view of $Q$ it makes sense to check whether $R$ is also close to $P$.

Candidates from all over the networks should be examinated. The problem with examining only neighbors' neighbors is that we will be eventually searching only within a single cluster, altough a good neighbor may have just joined at a random point in the network. Similar to the special long-links in small-world networks we need to establish links to other clusters. Likewise, when new nodes joins the network or they quickly changes their profile they should easily find an appropriate cluster to join. These issues call for a randomization when selectiong nodes to inspect for adding to a local view.

Node that triangle inequality property of the ranking function does not constitute an hard requirement for the system. In its absence closer neighbors are discovered based only on random encounters.

\paragraph{Proposed solutions} There has been some proposal for creating and managing a network topology in a large scale, dynamic, fully distributed system. Because these settings all the proposed solution are based on a gossip-based probabilistic solution to this problem

Topology management can be identified as an abstract middleware service that is desirable in large distributed system. This ease the development and of applications that makes usage of that service decoupling the decoupling typical distributed concerns.

We focus on solution that use a procative approach, i.e. autonomously build a topology my making usage of the peer profile and ranking fucntion as stated in the topology construction problem.

\textsc{T-Man} \cite{JB04} by using only local gossip messages gradually evolves the topology to match it to one defined by the ranking function. However the solution proposed does not take into account the concern about the exploration of the changes happening somewhere in the network because there is no support of randomization in order to probe the network for changes.

\cite{VS05, VSI06} decouple those aspects by adopting a two-layered set of gossip protocols for managing semantic overlay networks. They showed through simulations that this multi-layer approach ensures rapid convergence of network topology and at the same time guarantees highly adaptivity and low overhead.

These ideas have been applied to the design of BuddyCast \cite{PYM+08}, i.e. the first large scale internet deployed epidemic protocol stack. The main features are discovery of peers, discovery of content and formation on semantic overlay. All of these services are performed in a complete decentralized fashion. BuddyCast forms the core of the peer-to-peer file sharing system \textsc{Tribler} \cite{PGW+08}.

The lower layer, called Cyclon is responsible for maintaining connected overlay and for periodically feeding the top-layer protocol with nodes uniformly randomly selected from the network. Cyclon creates an overlay with completely random, uncorrelated links between nodes, such that the in-degree (number of incoming links) is practically the same for each node. Importantly it can achieve this property fairly quickly even when a small number of items (such 3 or 4) is exchanged in each communication. Therefore it is ideal as a lightweight service that can offer a node randomly selected peer from the current set of nodes. It offers a fully decentralized service for delivering information of new events that can happens everywhere in the network

The top layer protocol Vicinity is in charge of focusing on discovering peers that are semantically as close as possible and of adding these nodes to the local view.

\begin{algorithm}[t]
\caption{Cyclon and Vicinity epidepic protocols skeleton (cf. \cite{VS05, VSI06})}
\begin{algorithmic}[1]
\Procedure{ActiveThread}{}\Comment{Runs periodically every $T$ time units}

\State q $\gets$ \textsc{selectPeer}()
\State myItem $\gets$ (myAddress, timestamp, myProfile)
\State bufSend $\gets$ \textsc{selectItemsToSend}()
\State \textsc{send} bufSend to q
\State \textsc{receive} bufRecv from q
\State view $\gets$ \textsc{selectItemsToKeep}()

\EndProcedure

\Statex

\Procedure{PassiveThread}{}\Comment{Runs when contacted by some peer $p$}

\State \textsc{receive} bufRecv from p
\State myItem $\gets$ (myAddress, timestamp, myProfile)
\State bufSend $\gets$ \textsc{selectItemsToSend}()
\State \textsc{send} bufSend to p
\State view $\gets$ \textsc{selectItemsToKeep}()

\EndProcedure

\end{algorithmic}
\end{algorithm}

The difference between the two protocols vicinity and proximity reside in the different policy chosen for the function \texttt{selectItemsToSend()} and \texttt{selectItemsToKeep()}.

\paragraph{Gossip items} All information exchange between peers is carried out by means of gossip items. A gossip item created by peer $P$ is a tuple containing the following three fields:
\begin{itemize}
\item peer $P$'s contact information (e.g. network address and port)
\item item's creation time, i.e. timestamp
\item peer's profile
\end{itemize}
Each node maintains locally a number of items per protocol called the protocol's view.

\paragraph{Cache size} A large cache size provides higher chanches of making better items selection and therefore accelerate the construction of near-optimal semantic views. On the other hand the larger the cache size, the longer it takes to contact all peers in it, resulting in the existence of older and therefore more likely to be invalid links. \cite{VS05, VSI06} chose for their experiments a cache size of 50 items for each of the two layers.

\paragraph{Gossip length} The number of items exchanged in each communication is predefined and is called the protocol gossip length. It is $g_V$ for vicinity and $g_C$ for cyclon. It is the number of items gossiped per gossip exchange per protoco and it is a crucial factor for the amount of bandwidth used. So even though exchanging more items per gossip exchange allows information to disseminate faster, we are inclined to keep the gossip length for each of the two layer equal to 3.

\paragraph{Gossip period} The gossip period is a parameter that does not affect the protocol behaviour. The protocol evolves as a function of the number of messages exchanged. The gossip period only affects how fast the protocol's evolution will take place in time. The single constraint is that the gossip period should be adequately longer than the worse latency throughout the network, so that the gossip exchanges are not favored or hindered due to latency eterogeneity. A typical gossip period for our protocol would be 1 minute.

\paragraph{Node profile}

\paragraph{Ranking function}

\paragraph{Bandwidth considerations} Due to the periodic behaviour of gossiping, the price of having rapidly converging protocols may inhibit a high usage of network resources (i.e. bandwidth). In each cycle, a node gossips on average twice (i.e. once as an initiator and once as a responder). In each gossip $2(g_v+g_c)$ items are transferred to and from the node, resulting in a total traffin of $4(g_v+g_c)$ items for a node per cycle. An item's size is dominated by the profile it carries. [...] therefore a node's profile takes on average x bytes. So in each cycle the total number of bytes transferred to and from the node is $4 x (g_v+g_c)$.

For $g_v=g_c=3$ the average amount of data transferred to and from a node in one cycle is $x_2$ bytes, while for $g_v=g_c=1$ it is just $x_3$. Considering the gossip period T equal to 1 minute, this translates to an average bandwidth of $z_1$ and $z_2$ bytes per second respectively. With $g_v=g_c=3$ the systems adapts a little faster to changes, but if bandwidth is of high concern (minimal intrusion) $g_v=g_c=1$ can also provide good results.



